section{Conceptos ambientales}
    \begin{description}
        \item[Biósfera] Sistema formado por el conjunto de los seres vivos del planeta Tierra y sus interrelaciones.
        \item[Huella de carbono:] Medida del impacto de las actividades humanas en el medio ambiente, expresada en toneladas de $\ce{CO2}$ equivalente. %Un graficador psicrométrico preciso permite optimizar el consumo de combustible y electricidad en procesos industriales (secado, refrigeración), reduciendo directamente las emisiones asociadas a la generación de energía.
        \item[Ecología:] Ciencia que estudia las interrelaciones entre los factores bióticos y abióticos. %Desde la perspectiva de este proyecto, la ecología industrial permite ver a la planta de procesos como un organismo que debe interactuar eficientemente con su atmósfera local.
        \item[Factor Biótico:] Elementos vivos de un ecosistema que interactúan entre sí, como la flora, fauna y microorganismos. %El control psicrométrico preciso es vital para preservar estos factores en aplicaciones como la conservación de alimentos, invernaderos automatizados o laboratorios biotecnológicos.
        \item[Factor Abiótico:] Componentes físicos y químicos del ecosistema que carecen de vida, tales como la temperatura, la presión atmosférica y la humedad. %Estas son precisamente las variables de estado que el software procesa para definir el equilibrio termodinámico de las mezclas binarias.
        \item[Climatología:] Ciencia que estudia el clima y sus variaciones a largo plazo en una región determinada. %Los datos climatológicos históricos sirven como parámetros de diseño para el graficador, permitiendo establecer las condiciones promedio de operación de un sistema psicrométrico en una ubicación geográfica específica.
        \item[Meteorología:] Disciplina que estudia el estado del tiempo y los fenómenos atmosféricos en el corto plazo. %El graficador digital permite al ingeniero reaccionar ante cambios meteorológicos inmediatos, ajustando las trayectorias de los procesos térmicos para mantener la eficiencia ante variaciones de humedad y temperatura ambiente.
        \item[Cambio Climático:] Variación global del clima de la Tierra predominantemente debido a la acción humana. Este fenómeno altera las condiciones de diseño estándar, obligando al uso de herramientas digitales que puedan simular escenarios de alta temperatura y humedad extrema fuera de los rangos históricos.
        \item[Objetivos del Desarrollo Sostenible (ODS):] Conjunto de 17 metas globales de la ONU para erradicar la pobreza y proteger el planeta. Este proyecto se alinea específicamente con el ODS 7 (Energía asequible y no contaminante), el ODS 9 (Industria, innovación e infraestructura) y el ODS 12 (Producción y consumo responsables) al optimizar el uso de recursos energéticos.ono (ODP):] Cantidad de degradación química que un compuesto causa a la capa de ozono estratosférica. El software permite modelar sistemas de intercambio de calor que utilizan sustancias con $\text{ODP} = 0$, alineándose con el Protocolo de Montreal y promoviendo una ingeniería de procesos libre de sustancias agotadoras de la capa de ozono (SAO).
        \item[Resiliencia Climática:] Capacidad de un sistema absorber perturbaciones y recuperarse ante eventos climáticos extremos. Dado que los diagramas psicrométricos tradicionales son estáticos, este graficador digital permite simular condiciones de diseño "fuera de norma" (altas temperaturas o humedades extremas causadas por el calentamiento global), asegurando que los sistemas de climatización y procesos industriales sigan siendo operativos y seguros en el futuro.
        \item[Economía Circular:] Modelo de producción que busca minimizar el desperdicio. El uso de este software facilita la recuperación de energía térmica en corrientes de aire de desecho (calor sensible y latente). Al calcular con precisión el punto de rocío y la entalpía de mezcla, el ingeniero puede diseñar sistemas de recirculación y recuperación de condensado, cerrando el ciclo del agua y del calor en la planta industrial.
        \item[Eficiencia Energética:]   
        \item[Eficiencia Exergética:] Medida de la calidad de la energía y su degradación durante un proceso. A diferencia de la eficiencia energética simple, la exergía permite identificar dónde se pierde realmente la capacidad de realizar trabajo en una mezcla psicrométrica. %El graficador proporciona la base de datos necesaria para realizar análisis exergéticos avanzados, optimizando no solo cuánto combustible se quema, sino qué tan bien se utiliza la energía disponible.

        \item[Externalidades Ambientales:] Costos o beneficios indirectos que un proceso industrial impone a la sociedad y al medio ambiente. Un cálculo psicrométrico impreciso genera un sobredimensionamiento de equipos, lo que implica una mayor extracción de materias primas (metales, plásticos) y una mayor generación de residuos al final de su vida útil. Este proyecto mitiga estas externalidades mediante la precisión en el diseño térmico.
    \end{description}

\section{Definiciones químicas}
    \begin{description}
        \item[Sustancia:] Una sustancia es una forma de materia que tiene composición definida  que lo distinguen de los demás; es decir, difieren entre ellas por su composición, aspecto, color, sabor y otras propiedades.
        \item[Mezcla:] Es una combinación de dos o más sustancias, donde cada uno ellos conserva sus propiedades que los hacen únicos y los distinguen de los demás. Las mezclas pueden ser \textit{homogéneas}, cuando la composición de la mezcla es uniforme, o \textit{heterogéneas} cuando la mezcla resulta no uniforme; ejemplo: una mezcla de arena y aserrín. 
        \item[Elemento:] Un elemento es una sustancia que no puede separarse en otras más sencillas por medios químicos; ejemplo: los 118 elementos de la tabla periódica; pueden ser naturales o artificiales.
        \item[Compuesto:] Es una sustancia formada por dos átomos de dos o más elementos unidos químicamente en proporciones fijas y solo pueden separarse a través de medios químicos.
        \item[Estado de agregación:] Es la forma física en la que se presenta la materia basándose en la intensidad de las fuerzas de cohesión entre sus partículas (átomos, moléculas o iones) que dependen de su energía cinética (asociado a la temperatura) y potencial (atracción). Actualmente se conocen 7 estados de la materia; para el desarrollo de este trabajo solo se consideran los estados: sólido líquido, gaseoso.
        \item[Fase:] Región de un sistema termodinámico donde todas las propiedades físicas son uniformes; es decir, son homogéneos, uniforme, y presentan ausencia de interfaces.
        \item[Vapor:] Es la sustancia condensable que, bajo condiciones de presión y temperatura, tiene la capacidad de migrar de la fase gaseosa a la fase líquida o sólida. Su comportamiento está limitado por su presión de saturación, la cual es una función de la temperatura
        \item[Gas:] Es el componente o mezcla de ellas que no se condensa en todo rango de operación y es químicamente inerte con el vapor, lo cual permite la aplicación de la ley de Dalton. 
        \item[Peso Atómico:] 
        \item[Peso Molecular:] 
        \item[Equilibrio:] 
        \item[Equilibrio mecánico (física):]
        \item[Equilibrio térmico (termodinámica):]
        \item[Equilibrio químico (química):]
        \item[Equilibrio de fases:]
    \end{description}

\subsection{Conceptos de Equilibrio de Fases}
    \begin{description}
        \item[Presión (${\pressure[]{}}$):] 
        \item[Presión de vapor ($\pVapor$):] La presión parcial ejercida por las moléculas del componente condensable (vapor).
        \item[Presión de Saturación $(\pSat)$:] Se define como la presión cuando a una temperatura dada coexisten la fase líquida y vapor en equilibrio termodinámico. 
        \item[Estado de Saturación:] Condición en la que el componente condensable de la mezcla binaria se encuentra en equilibrio termodinámico entre sus fases líquida y vapor. En el diagrama, esto define la Curva de Saturación $\relativeHumidity[]{} = 100$. 
        \item[Punto triple (${}$):]
        \item[Punto crítico (${}$):]
        \item[Curva de condensación (${}$):]
        \item[Curva de sublimación (${}$):]
        \item[ (${}$):]
        \item[ (${}$):]
        \item[ (${}$):]
    \end{description}



\section{Magnitudes Termodinámicas Fundamentales}
    \begin{description}
        \item[Calor:] 
        \item[Temperatura:] 
        \item[Temperatura crítica:] 
        \item[Punto triple:] Es el punto donde coexisten las fases líquida, sólida y vapor simultáneamente. 
        \item[Energía (${}$):]
        \item[Energía Interna (${}$):]
        \item[Entalpía (${}$):]
        \item[Capacidad calorífica (${}$):] La capacidad calorífica de una sustancia, es la cantidad de calor que se necesita añadirse a una cantidad de sustancia en $1 \unit{\celsius}$ 
        \item[calor específico (${}$):] Representa la capacidad de la mezcla (gas seco + componente condensable) para almacenar energía térmica ante un cambio de temperatura a presión constante. El software utiliza polinomios dependientes de la temperatura para ajustar este valor dinámicamente según el componente seleccionado (ej. $\ce{NH3}$ o $\ce{H2O}$).
        \item[Calor sensible (${}$):] El calor sensible es la energía que se asocia al cambio de temperatura de una sustancia sin que suceda algún cambio de estado; se define así:
        \item[Calor latente (${}$):]
        \item[Entalpía (${}$):] 
        \item[Entalpía específica (${}$):] 
        \item[Calor de vaporización (${}$):] 
    \end{description}


    
\section{Variables de Estado Psicrométrico}
    \begin{description}
        \item[Humedad absoluta] Según \parencite[p. 729]{bib:cengel2015termodinamica} se le conoce también como humedad específica\footnote{para \parencite{bib:green2007perry} esta expresión en específico tiene su propia definición y distinto que se verá más adelante.}, contenido de humedad o \textit{índice de humedad} y para \parencite[p. 12-4]{bib:green2007perry} como \textit{relación de masa} o \textit{relación de mezcla}; para el desarrollo de este proyecto se ha optado usar las definiciones de este último. La humedad absoluta se representa la relación matemática de la masa de vapor respecto a una unidad de masa del gas; generalmente se expresa en, $\specificHumidityUnit$.
        \eqAbsoluteHumidity*
        \item[Volumen específico:] Es la cantidad de volumen que ocupa el gas una cantidad de gas. A medida que la temperatura aumenta también aumenta el volumen específico.
        \item[Humedad específica:] Conocido también como fracción másica y se define como la masa del vapor por unidad de masa de la mezcla vapor-gas.
        \item[Humedad relativa:] Es la relación de la presión de vapor del sistema respecto a la presión de vapor saturado del sistema bajo las mismas condiciones de temperatura. Para una mezcla binaria de vapor de agua y aire seco, la capacidad de retención de vapor aumenta con la temperatura.
        \item[Temperatura de Bulbo Seco:]  
        \item[Temperatura de Bulbo Húmedo:] Corresponde a la temperatura de equilibrio dinámico alcanzado por una superficie líquida de la cual el component vapor que se encuentra en estado líquido se evapora en una corriente de gas seco que fluye cuando la tasa de transferencia de calor a la superficie por convección es igual a la velocidad de transferencia de masa lejos de la superficie. Esta temperatura está muy cerca a la temperatura a la temperatura de saturación adiabática para la mezcla del sistema agua-aire; pero, no para otras mezclas.
        \item[Humedad absoluta molar:] Se define como la relación de moles del vapor respecto a moles del gas seco; matemáticamente se expresa de la siguiente manera
        \item[Grado de saturación:] Se define como la relación de la humedad absoluta respecto a la humedad absoluta saturada,
        
        La diferencia con la humedad relativa (\relativeHumidity) es, que el grado de saturación (\degreeOfSaturation[]{}) es una relación de masas; la humedad relativa es una relación de presiones.
        \item[Humedad específica:] Es la relación  entre el contenido de vpaor de la mezcla respecto al contenido total del aire. Algunas veces también se le llama como relación de humedad o contenido de humedad o relación de mezcla
        \item[Punto de rocío:] También se le conoce com la \textbf{temperatura de saturación}; es la temperatura al cual se satura el vapor al enfriarse; es decir, la temperatura donde el vapor ejerce una presión de vapor igual a la presión de saturación; la condición de saturación ocurre cuando:

    \end{description}


\section{Fundamentos de Ingeniería de Software}
    \begin{description}
        \item[Abstracción de Datos:] Representación de estados termodinámicos como objetos.
        \item[Interfaz de Programación de Aplicaciones (API):] Definición de métodos de entrada y salida para el usuario.
        \item[Visualización de Datos:] Conceptos de renderizado de gráficos vectoriales.
        \item[Diagrama de Clases:] Tipo de diagrama \gls{uml} que describe la estructura del sistema mostrando sus clases, atributos y las relaciones entre objetos. Se utiliza para modelar las entidades termodinámicas (ej. Clase \texttt{MezclaBinaria}, Clase \texttt{Componente}, Clase \texttt{Estado}) y sus dependencias, como la relación entre un componente y sus coeficientes de Antoine.
        \item[Diagrama de Secuencia:] Diagrama que ilustra cómo los objetos interactúan entre sí en un orden temporal. Es crítico para documentar el flujo de datos: desde que el usuario mueve un selector en el frontend (Flutter), la petición viaja al backend (Rust), se ejecutan las iteraciones de Newton-Raphson, y el resultado vuelve para redibujar el diagrama psicrométrico.
        \item[Principios SOLID:] Acrónimo de cinco principios de diseño orientado a objetos destinados a hacer que el software sea más comprensible, flexible y mantenible:
            \begin{itemize}
                \item \textbf{S (Single Responsibility):} Cada módulo (crate) tiene una sola razón para cambiar (ej. el módulo de cálculo de presión de saturación solo hace eso).
                \item \textbf{O (Open/Closed):} El software permite añadir nuevas mezclas binarias (ej. Metanol-Aire) sin modificar el motor de cálculo central.
                \item \textbf{L (Liskov Substitution):} Las diferentes \glspl{eos} (Virial, Peng-Robinson) deben ser intercambiables sin romper la lógica del graficador.
                \item \textbf{I (Interface Segregation):} El backend solo expone los datos necesarios para el gráfico, no toda su lógica interna.
                \item \textbf{D (Dependency Inversion):} La interfaz no depende de una implementación específica de cálculo, sino de una abstracción (trait en Rust).
            \end{itemize}
        \item[Bases de Datos:] Sistema de almacenamiento estructurado (puede ser JSON, SQL, XML, ton o archivos embebidos en el binario) que contiene los parámetros críticos ($\Tc, \pC$), factores acéntricos y constantes empíricas de cientos de sustancias. El graficador consulta esta base para configurar el sistema binario seleccionado por el usuario.
        \item[Patrón de Diseño "Worker" o Isolate:] Estrategia para evitar el bloqueo de la interfaz de usuario. Dado que el renderizado psicrométrico exige cálculos intensivos, se delegan estas tareas a hilos secundarios (en Rust o en \textit{Isolates} de Dart), permitiendo que el usuario siga interactuando con la aplicación mientras el gráfico se genera en segundo plano.
        \item[Programación Orientada a Objetos (POO) y Composición:] Aunque Rust no es un lenguaje de POO tradicional, utiliza \textit{Structs} para definir datos y \textit{Traits} para definir comportamientos (polimorfismo). En el frontend con \textbf{Flutter}, la POO es fundamental: cada elemento de la interfaz (botones, ejes, líneas) es un \textit{Widget} que hereda propiedades de clases base, permitiendo una construcción jerárquica y eficiente del diagrama psicrométrico.
        \item[Mapeo de Coordenadas (Transformación Afín):] Es el proceso matemático de convertir valores físicos (como Temperatura en $^\circ$C y Humedad en kg/kg) a coordenadas de píxeles ($x, y$) en el lienzo de \textbf{Flutter}. Mediante transformaciones afines de escalado y traslación, el software garantiza que el diagrama sea interactivo, permitiendo funciones de \textit{zoom} y desplazamiento (\textit{panning}) sin perder la referencia física de los datos.
        \item[Backend y Frontend:] El \textit{Backend} (desarrollado en \textbf{Rust}) actúa como el núcleo de cómputo científico, procesando la física del sistema. El \textit{Frontend} (desarrollado en \textbf{Flutter}) funciona como la capa de presentación y experiencia de usuario. La separación clara entre ambos permite que el motor de cálculo sea multiplataforma y pueda integrarse incluso en servidores de alto rendimiento o sistemas embebidos.
        \item[Refactorización (\textit{Refactoring}):] Proceso de reestructurar el código fuente existente sin cambiar su comportamiento externo. En este proyecto, la refactorización constante permite optimizar los algoritmos de mezcla para que el graficador sea cada vez más rápido y eficiente.
        \item[Control de Versiones (Git):] Sistema para rastrear cambios en el código fuente durante el desarrollo. Facilita la experimentación con diferentes modelos termodinámicos (creando ramas o \textit{branches}) y asegura que siempre exista un historial de las ecuaciones implementadas.
        \item[Pruebas Unitarias (\textit{Unit Testing}):] Procedimiento para validar que cada pequeña parte del código (ej. una función que calcula la entalpía) funcione correctamente de forma aislada. Es la garantía de que los resultados del graficador coinciden con los datos experimentales de referencia.
        \item[Experiencia de Usuario (UX) e Interfaz de Usuario (UI):] Conceptos de diseño que buscan que el graficador no solo sea preciso, sino fácil de usar. Esto incluye la elección de colores para las isolíneas, el manejo del \textit{zoom} táctil y la legibilidad de las etiquetas de datos en el diagrama.
        \item[Gestión de Dependencias y Ecosistema (Cargo):] En Rust, el gestor de paquetes \texttt{Cargo} permite la integración de \textit{Crates} especializadas en álgebra lineal y termodinámica. Esto asegura la reproducibilidad del software, permitiendo que cualquier investigador pueda compilar el proyecto con las mismas versiones exactas de las bibliotecas de cálculo, eliminando el problema de "funciona en mi máquina".

        \item[Compilación Ahead-of-Time (AOT) vs. Just-in-Time (JIT):] \textbf{Flutter} utiliza compilación AOT para generar código de máquina nativo, lo que es esencial para el renderizado de miles de puntos en el diagrama psicrométrico a 60 cuadros por segundo (FPS). Mientras tanto, el backend en Rust se compila de forma nativa para extraer el máximo rendimiento de la CPU en el procesamiento de modelos viriales complejos.
        \item[Transformación Afín y Mapeo No Lineal:] Técnica geométrica que preserva líneas y proporciones. En el graficador, se utiliza para mapear el dominio termodinámico $(T, \omega)$ al dominio de la pantalla $(x, y)$. Dado que algunos diagramas usan escalas logarítmicas (como el diagrama de Mollier), el software implementa transformaciones no lineales para proyectar curvas de saturación con precisión sub-píxel.
        \item[Patrón de Diseño "State Management" (Gestión de Estado):] Estrategia para sincronizar los datos termodinámicos con la interfaz de usuario. En Flutter, patrones como \textit{Provider} o \textit{BLoC} aseguran que si la presión del sistema cambia, todos los observadores (la curva de saturación, las isolíneas y las etiquetas de datos) se redibujen de manera consistente y eficiente.

        % \item[Redes Neuronales de Regresión (ML):] Subcampo del \gls{ml} orientado a predecir valores continuos. Para sistemas binarios complejos con comportamientos no ideales extremos, el graficador puede emplear modelos de regresión para aproximar el coeficiente de actividad ($\gamma$) o el factor de compresibilidad ($Z$), sirviendo como una alternativa ultrarrápida a las ecuaciones cúbicas de estado.
        % \item[Algoritmos de Optimización y Root-Finding:] Métodos numéricos para hallar raíces de funciones. El software implementa algoritmos como \textbf{Brent} o \textbf{Newton-Raphson} para calcular la temperatura de bulbo húmedo y el punto de rocío, donde no existen soluciones analíticas directas y se requiere una convergencia de alta precisión ($10^{-6}$ o menor).
        % \item[Machine Learning (ML) en Ingeniería de Procesos:] El Aprendizaje Automático se define como el uso de algoritmos para identificar patrones en datos históricos. En este proyecto, el \gls{ml} tiene el potencial de actuar como un "Metamodelo": en lugar de resolver costosas ecuaciones iterativas de gases reales (como Lee-Kesler) miles de veces por segundo, se puede entrenar una \textbf{Red Neuronal Regresora} para predecir propiedades psicrométricas con un error mínimo y un costo computacional casi nulo.
        % \item[Inteligencia Artificial (IA) para la Optimización:] La \gls{ia} abarca técnicas que permiten a la máquina tomar decisiones óptimas. En el contexto del graficador, la IA puede utilizarse para la detección automática de errores en los datos de entrada o para sugerir el modelo termodinámico más adecuado (Virial vs. Cúbico) basándose en las condiciones de presión y temperatura detectadas, optimizando así la precisión del software de forma autónoma.
        % \item[Serialización y Deserialización (Serde):] Proceso de convertir estructuras de datos en memoria a formatos de intercambio como JSON o Protocol Buffers. Esto permite que el backend (Rust) envíe los resultados de las trayectorias termodinámicas al frontend (Flutter) de manera compacta, facilitando incluso la comunicación vía WebSockets para futuras versiones basadas en la nube.

        \item[Pipelines de Integración Continua (CI/CD):] Automatización de pruebas y despliegue. En un proyecto de grado profesional, esto garantiza que cada cambio en las ecuaciones de entalpía pase por una batería de pruebas unitarias (\textit{unit tests}) automáticas que validen los resultados contra datos experimentales conocidos.
        \item[Gráficos Vectoriales y Renderizado por GPU:] A diferencia de los gráficos basados en píxeles (raster), el graficador utiliza comandos de dibujo vectorial. Esto permite que las líneas de entalpía constante se definan matemáticamente y sean procesadas por la GPU mediante el motor \textit{Impeller} o \textit{Skia}, permitiendo un escalado infinito sin pérdida de definición.
    \end{description}


%     \subsection{Punto de Rocío} %  (Dew Point)
%         \eqnSaturationCondition
%         \figurePuntoRocio{fig:rocio_isobarico}
%         Representa la temperatura más baja que puede alcanzarse por enfriamiento evaporativo de una superficie ventilada humedecida con el vapor. 
%         En condiciones donde la presión de vapor es baja y la temperatura desciende bajo el punto triple del condensable, esta transición implica la sublimación del sólido (formación de escarcha) para alcanzar el equilibrio de saturación.


%         \begin{figure}[h]
%             \centering
%             \def\Tamb{85}       % Temperatura ambiente: 85
\def\Pv{4}          % Presión de vapor actual: 40

\def\Tas{60.0}  
\def\Pas{20}

% --- water_saturation_adiabatic.tex ---
\begin{tikzpicture}
\begin{axis}[
    width=11cm, height=8.5cm,
    ymode=log, 
    log ticks with fixed point,
    title={\textbf{Proceso de Saturación Adiabática ($H_2O$)}},
    xlabel={$\temperature[]{}$ [$^\circ$C]},
    ylabel={$\pressure[]{\vapor}$ [kPa] (Log)},
    xmin=\Tmin, xmax=\Tmax, 
    ymin=\Pmin, ymax=\Pmax,
    axis lines=left,
    % Ticks específicos para resaltar Tas y Tamb
    xtick={0, \Tas, \Tamb}, 
    xticklabels={0, $\temperature[]{sat}$, $\temperature[]{amb}$},
    ytick={0.1, \Ptp, 10, \Pv, \Pas, 101.3}, 
    yticklabels={0.1, \Ptp, 10, $\pressure[]{\vapor, 1}$, $\pressure[]{\text{sat}}(\temperature[]{sat})$, 101.3},
    grid=both,
    minor grid style={gray!5},
    major grid style={gray!15},
    clip=true
]

    % --- CURVAS DE COEXISTENCIA ---
    % Sublimación
    \addplot[teal, ultra thick, domain=\Tmin:\Ttp, samples=100] {0.611*exp(22.5*(1 - 273.16/(x+273.15)))};
    
    % Fusión
    \draw[blue, thick] (axis cs:\Ttp,\Ptp) -- (axis cs:-1.5,\Pmax); 
    
    % Vaporización (Saturación)
    \addplot[blue, ultra thick, domain=\Ttp:\Tmax, samples=100] {0.611*exp(17.27*x/(x+237.3))}
        node[pos=0.3, sloped, above, black, font=\tiny] {Curva de saturación, $\pSat$};

    % --- ETIQUETAS DE ESTADO ---
    \node[blue!20!black, font=\footnotesize] at (axis cs:-15, 50) {\textbf{SÓLIDO}};
    \node[blue!70!black, font=\footnotesize] at (axis cs:35, 80) {\textbf{LÍQUIDO}};
    \node[orange!70!black, font=\footnotesize] at (axis cs:45, 0.5) {\textbf{VAPOR}};

    % --- PROCESO DE SATURACIÓN ADIABÁTICA ---
    % Trayectoria del estado 1 (ambiente) al estado 2 (saturado adiabáticamente)
    \draw[red, ultra thick, -{Stealth[scale=1.2]}] (axis cs:\Tamb,\Pv) -- (axis cs:\Tas,\Pas)
        node[midway, above, sloped, font=\tiny, text=red!50!black] {Sat. Adiabática};
    
    % Proyecciones de los puntos de interés
    % \draw[gray, dashed] (axis cs:\Tas,\Pas) -- (axis cs:\Tas,\Pmin);
    % \draw[gray, dashed] (axis cs:\Tamb,\Pv) -- (axis cs:\Tamb,\Pmin);
    % \draw[gray, dashed] (axis cs:\Tas,\Pas) -- (axis cs:\Tmin,\Pas);
    % \draw[gray, dotted] (axis cs:\Tamb,\Pv) -- (axis cs:\Tmin,\Pv);

    % Punto 1: Aire de entrada (Estado inicial)
    \node[circle, fill=black, inner sep=1.5pt, label=right:{\scriptsize 1}] at (axis cs:\Tamb,\Pv) {};
    
    % Punto 2: Aire saturado a la salida (Estado final)
    \node[circle, fill=red, inner sep=1.5pt, label=above left:{\scriptsize 2}] at (axis cs:\Tas,\Pas) {};

    % Punto Triple
    \filldraw (axis cs:\Ttp,\Ptp) circle (1.5pt) node[right, font=\tiny, xshift=2pt] {\triplePoint};

\end{axis}
\end{tikzpicture}
 % Llamada al segundo módulo
%             \caption{Localización del punto de rocío en la línea de vapor saturado ($x=1$).}
%         \end{figure}
        
% \eqnDegreeOfSaturation
%         \eqnVolumetricHumidity

%     \figureAdiabaticSaturationTS{fig:water_TS}
    
%         \eqnRelativeHumidity



